\documentstyle[%


righttag%


,11pt%


,amssymb%


,russian%               remove in Eng.


,izvru%                 Set izven% in Eng.


]{amsart}





% !{=}





\newcommand{\co}{\operatorname{co}}


\newcommand{\la}{\langle}


\newcommand{\ra}{\rangle}


\newcommand{\tts}[1]{}





\theoremstyle{plain}


\theorembodyfont{\it}


\newtheorem{assrt}{\hskip\parindent Утверждение}





\renewcommand{\baselinestretch}{0.97}





%\hoffset=-15mm%                  For Eng.


\setcounter{page}{37}


\god{2004}


\nomer{12~(511)} %                        Remove in Eng.


%\nomer{Vol.\,48, No.\,12}%               Set in Eng.


%\str{\pageref{begin}--\pageref{end}}%  Set in Eng.


\udk{519.10}





\author{\it В.М.\,Кравцов}


% NB:   ^^ remove in Eng.


\title{О новых свойствах максимально нецелочисленных вершин многогранника 


трехиндексной аксиальной задачи о назначениях}





\date{}


%\pagestyle{myheadings}%         Set in Eng.


\pagestyle{plain} %              Remove in Eng.


\begin{document}


%\thispagestyle{plain}%          Set in Eng.


\maketitle


\label{begin}








\section*{Введение}





Известно $[1]$, что трехиндексная аксиальная задача о назначениях,


имеющая многочисленные приложения [2]--[4], является NP-полной.


Сложное строение нецелочисленных вершин многогранника


$M(3,n)=\Big\{x=\|x_{ijt}\|_n: x_{ijt}


\geq 0$ $\forall (i,j,t)\in N_n^3$, $\sum\limits_{i=1}^n\sum\limits_{j=1}^n


x_{ijt}=1$ $\forall t\in N_n$, $\sum\limits_{i=1}^n\sum\limits_{t=1}^n


x_{ijt}=1$ $\forall j\in N_n$, $\sum\limits_{j=1}^n\sum\limits_{t=1}^n


x_{ijt}=1$ $\forall i\in N_n\Big\}$, где


$N_n=\{1,2,\dotsc,n\}$, $n\geq 2$, $N_n^3 = N_n\times N_n\times


N_n$, порожденного условиями этой задачи, создает


принципиальные трудности при его исследовании.





Для современных исследований свойств


многогранника $M(3,n)$ характерны следующие два направления. К


первому напрaвлению относятся работы [5]--[7], в которых изучаются


граневые структуры и фасеты (грани максимальной


размерности) многогранника $M(3,n)$. Основным объектом


исследований второго направления являются его нецелочисленные


вершины [8]--[14]. Примеры нецелочисленных вершин этого


многогранника встречались давно [15], но само


понятие $s$-нецелочисленной вершины введено совсем недавно.


Напомним [9], [11] это понятие. Вершина многогранника


$M(3,n)$ называется $s$-нецелочисленной, если она содержит ровно


$s$ нецелочисленных (дробных) компонент. Упомянем два результата,


относящихся ко второму направлению.


\begin{itemize}


\item[1)] Для любого числа $s\in \{4,6,7,\dotsc,3n-2\}$, и


только для него, у многогранника $M(3,n)$ существуют


$s$-нецелочисленные вершины [12].


\item[2)] Оценки снизу числа $\sigma (n,s)$ \ $s$-нецелочисленных


вершин многогранника $M(3,n)$ [9], [11], [13], позволившие


опровергнуть гипотезу 18 из [16]. С помощью этих оценок и


явных формул для $\sigma (n,4)$ и $\sigma (n,6)$ [11] в [9]


получены оценки снизу числа нецелочисленных вершин многогранника


$M(p,n)$ \ $p$-индексной ($p\geq 3$) аксиальной задачи о


назначениях, которые значительно улучшают оценку из [8].


\end{itemize}





Так как всякая вершина многогранника $M(3,n)$ содержит не более


чем $3n-2$ положительных компонент [15], то


$(3n-2)$-нецелочисленные вершины этого многогранника будем


называть максимально нецелочисленными вершинами (м.\,н.\,в.).


К настоящему времени известно [10], [14]


\ $(n-\left[\frac{n}{2}\right]-3)(\left[\frac{n}{2}\right]-2)


+n+2\left[\frac{n}{2}\right]-8$


типов неэквивалентных м.\,н.\,в. многогранника $M(3,n)$, $n\geq 7$.


Идентификация типов вершин проводится по количеству дробных


компонент, содержащихся в двумерных сечениях трехиндексных матриц,


представляющих собой его вершины. Две вершины многогранника


$M(3,n)$ называются неэквивалентными, если одна из них не может


быть переведена в другую путем перестановки ее двумерных сечений.


Неэквивалентные вершины этого многогранника имеют различные


структуры.





В данной работе указаны новые типы м.\,н.\,в. многогранника


$M(3,n)$ и для них установлено существование различных структур.


Доказано, что у многогранника $M(3,n)$, $n\geq 7$, существуют


м.\,н.\,в., содержащие компоненты, равные $\frac{1}{n^2 -5n+6}$ и


$\frac{n^2 -7n+10}{n^2 -7n+12}$. Тем самым получено опровержение


предположения, высказанного в [10], согласно которому для любого


натурального числа $n$ наименьшая (наибольшая) положительная


компонента среди всех м.\,н.\,в. многогранника $M(3,n)$, $n\geq 3$,


равна $\frac{1}{n}$ $\left(\frac{n-1}{n}\right)$. Разработана также


простая процедура построения м.\,н.\,в. многогранника $M(3,n)$, $n\geq


4$, на базе м.\,н.\,в. многогранника $M(3,n-1)$, которая находит


применение при доказательстве приведенных теорем. В дальнейшем


будем предполагать, что $n\geq 3.$








\section{Предварительные результаты}





Известно [15], что ранг матрицы системы ограничений


трехиндексной аксиальной транспортной задачи порядка $m\times


k\times l$ равен числу $m+k+l-2$. С помощью этого факта


методом от противного легко доказывается





\begin{assrt}


Для того чтобы матрица $x$ была м.\,н.\,в.


многогранника $M(3,n)$, необходимо, чтобы она содержала $3n-2$


положительных элементов и для любых непустых подмножеств $I$, $J$, $T$


множества $N_n$ выполнялось неравенство $|x(I,J,T)|\leq


|I|+|J|+|T|-2$, где $x(I,J,T)= \{(i,j,t)\in I\times J\times T:


x_{ijt}> 0 \}$ $(|\emptyset|=0)$.


\end{assrt}





Утверждение, обратное утверждению 1, неверно.





Зафиксируем число $m\in N_{n-1}$. Пусть $I_1$, $J_1$, $T_1$ --- некоторые


подмножества (возможно, совпадающие) мощности $m$ множества $N_n$.


Положим $I_2 = N_n \setminus I_1$, $J_2 = N_n \setminus J_1$, $T_2


= N_n \setminus T_1$. Так как $m\leq n-1$, то $I_2 \neq


\emptyset$, $J_2 \neq \emptyset$, $T_2 \neq \emptyset$. Для двух


троек подмножеств $(I_1,J_1,T_1)$ и $(I_2,J_2,T_2)$ определим


многогранники $M(I_s,J_s,T_s)= \Big\{x{=}\|x_{ijt}\|_{|I_s|\times


|J_s|\times |T_s|}: \sum\limits_{i\in I_s} \sum\limits_{j\in J_s} x_{ijt}=1$


$\forall t\in T_s$, $\sum\limits_{i\in I_s} \sum\limits_{t\in T_s} x_{ijt}=1$


$\forall j\in J_s$, $\sum\limits_{j\in J_s} \sum\limits_{t\in T_s} x_{ijt}=1$


$\forall i\in I_s$, $x_{ijt}\geq 0$ $\forall (i,j,t)\in I_s\times


J_s\times T_s \Big\}$, $s=1,2$.





\begin{note}


Многогранники $M(I_1,J_1,T_1)$ и $M(I_2,J_2,T_2)$


отличаются от многогранников $M(3,m)$ и


$M(3,n-m)$ соответственно лишь нумерацией элементов их матриц.


\end{note}





При $m=1$ многогранник $M(I_1,J_1,T_1)$, а при $m=n-1$


многогранник $M(I_2,J_2,T_2)$ вырождаются в точку.





Для вершины $y^s=\|y_{ijt}^s\|_{|I_s|\times |J_s|\times |T_s|}$


многогранника $M(I_s,J_s,T_s)$, $s=1,2$, введем множество


$$


K(I_s,J_s,T_s,y^s)=\{(i,j,t)\in I_s\times J_s\times T_s: y_{ijt}^s>0\}.


$$





С помощью утверждения 1 доказывается 





\begin{lem}


Пусть $y^s=\|y_{ijt}^s\|_{|I_s|\times |J_s|\times |T_s|}$


--- некоторая вершина многогранника $M(I_s,J_s,T_s)$, $s=1,2$. Тогда


для любых двух троек индексов $(i_1,j_1,t_1)\in


K(I_1,J_1,T_1,y^1)$ и $(i_2,j_2,t_2)\in K(I_2,J_2,T_2,y^2)$ и


любого числа $r\in \{1,2,3\}$ матрица $x^r=\|x^r_{ijt}\|_n$ с


элементами





$x_{ijt}^r=y_{ijt}^s$ $\forall (i,j,t)\in


K(I_s,J_s,T_s,y^s)\setminus \{(i_s,j_s,t_s)\}$, $s=1,2$,





$x_{ijt}^r=y_{ijt}^s-\theta$, если $(i,j,t)=(i_s,j_s,t_s)$,


$s=1,2$,





$x_{ijt}^1=\theta$ $\forall (i,j,t)\in


\{(i_1,j_2,t_1),(i_2,j_1,t_2)\}$,





$x_{ijt}^2=\theta$ $\forall (i,j,t)\in


\{(i_2,j_1,t_1),(i_1,j_2,t_2)\}$,





$x_{ijt}^3=\theta$ $\forall (i,j,t)\in


\{(i_1,j_1,t_2),(i_2,j_2,t_1)\}$,





$x_{ijt}^r=0$ для остальных $(i,j,t)\in N_n^3$





\noindent{}является вершиной многогранника $M(3,n)$, где $\theta=\min


\{y_{i_1,j_1,t_1}^1, y_{i_2,j_2,t_2}^2 \}$.


\end{lem}





\begin{cor}


Пусть $y^s=\|y_{ijt}^s\|_{|I_s|\times |J_s|\times


|T_s|}$ --- некоторая м.\,н.\,в. многогранника $M(I_s,J_s,T_s)$,


$s=1,2$. Если для некоторой пары троек индексов $(i_1,j_1,t_1)\in


K(I_1,J_1,T_1,y^1)$ и $(i_2,j_2,t_2)\in K(I_2,J_2,T_2,y^2)$


выполняется условие $y_{i_1,j_1,t_1}^1 \neq y_{i_2,j_2,t_2}^2$


($y_{i_1,j_1,t_1}^1 = y_{i_2,j_2,t_2}^2$), то вершины $x^1$,


$x^2$, $x^3$, указанные в лемме 1, являются


$(3n-3)$-нецелочисленными ($(3n-4)$-нецелочисленными) вершинами


многогранника $M(3,n)$.


\end{cor}





С помощью следствия 1 и утверждения 1 легко доказать следующую


лемму, дающую простую процедуру построения м.\,н.\,в. многогранника


$M(3,n)$ с использованием м.\,н.\,в. многогранников $M(I_1,J_1,T_1)$ и


$M(I_2,J_2,T_2)$.





\begin{lem}


Пусть $n\geq 4$, $|I_s|=|J_s|=|T_s|\geq 2$ и


$y^s=\|y_{ijt}^s\|_{|I_s|\times |J_s|\times |T_s|}$ --- некоторая


м.\,н.\,в. многогранника $M(I_s,J_s,T_s)$, $s=1,2$. Тогда для любых


четырех троек индексов $(i_s,j_s,t_s)$, $(i'_s,j'_s,t'_s)\in


K(I_s,J_s,T_s,y^s)$, $s=1,2$, удовлетворяющих условиям


\begin{equation}


\label{math/1} y_{i_1,j_1,t_1}^1 \neq y_{i_2,j_2,t_2}^2,\quad


y_{i'_1,j'_1,t'_1}^1 \neq y_{i'_2,j'_2,t'_2}^2,


\end{equation}


и любого числа $r\in N_6$ матрица $x^r=\|x^r_{ijt}\|_n$ с


элементами





$x_{ijt}^r=y_{ijt}^s$ $\forall (i,j,t)\in


K(I_s,J_s,T_s,y^s)\setminus \{(i_s,j_s,t_s)$, $(i'_s,j'_s,t'_s)\}$,


$s=1,2$,





$x_{ijt}^r=y_{ijt}^s-\theta$, если $(i,j,t)=(i_s,j_s,t_s)$,


$s=1,2$,





$x_{ijt}^r=y_{ijt}^s-\theta'$, если $(i,j,t)=(i'_s,j'_s,t'_s)$,


$s=1,2$,





$x_{ijt}^1=\theta$ $\forall (i,j,t)\in


\{(i_1,j_2,t_1),(i_2,j_1,t_2)\}$,





$x_{ijt}^1=\theta'$ $\forall (i,j,t)\in


\{(i'_2,j'_1,t'_1),(i'_1,j'_2,t'_2)\}$,





$x_{ijt}^2=\theta$ $\forall (i,j,t)\in


\{(i_1,j_2,t_1),(i_2,j_1,t_2)\}$,





$x_{ijt}^2=\theta'$ $\forall (i,j,t)\in


\{(i'_1,j'_1,t'_2),(i'_2,j'_2,t'_1)\}$,





$x_{ijt}^3=\theta$ $\forall (i,j,t)\in


\{(i_2,j_1,t_1),(i_1,j_2,t_2)\}$,





$x_{ijt}^3=\theta'$ $\forall (i,j,t)\in


\{(i'_1,j'_2,t'_1),(i'_2,j'_1,t'_2)\}$,





$x_{ijt}^4=\theta$ $\forall (i,j,t)\in


\{(i_2,j_1,t_1),(i_1,j_2,t_2)\}$,





$x_{ijt}^4=\theta'$ $\forall (i,j,t)\in


\{(i'_1,j'_1,t'_2),(i'_2,j'_2,t'_1)\}$,





$x_{ijt}^5=\theta$ $\forall (i,j,t)\in


\{(i_1,j_1,t_2),(i_2,j_2,t_1)\}$,





$x_{ijt}^5=\theta'$ $\forall (i,j,t)\in


\{(i'_1,j'_2,t'_1),(i'_2,j'_1,t'_2)\}$,





$x_{ijt}^6=\theta$ $\forall (i,j,t)\in


\{(i_1,j_1,t_2),(i_2,j_2,t_1)\}$,





$x_{ijt}^6=\theta'$ $\forall (i,j,t)\in


\{(i'_2,j'_1,t'_1),(i'_1,j'_2,t'_2)\}$,





$x_{ijt}^r=0$ для остальных $(i,j,t)\in N_n^3$





\noindent{}является м.\,н.\,в. многогранника $M(3,n)$, где $\theta=\min


\{y_{i_1,j_1,t_1}^1, y_{i_2,j_2,t_2}^2 \}$, $\theta'=\min


\{y_{i'_1,j'_1,t'_1}^1, y_{i'_2,j'_2,t'_2}^2 \}$.


\end{lem}





\begin{note}


Отметим, что лемма 2 остается справедливой и для


случая, когда $n\geq 4$, $|I_2|=|J_2|=|T_2|=1$. При этом


$(i_2,j_2,t_2)=(i'_2,j'_2,t'_2)$, условие


(1) заменяется на условие $y_{i_1,j_1,t_1}^1 +


y^1_{i'_1,j'_1,t'_1} \neq 1$ и величина $\theta'$ вычисляется по


формуле $\theta'=\min \{y_{i'_1,j'_1,t'_1}^1 , 1-\theta \}$.


\end{note}








\section{Основные результаты}





Совокупность элементов трехиндексной матрицы $x=\|x_{ijt}\|_n$ с


фиксированным значением одного индекса, например $t$, будем


называть двумерным сечением ориентации $(i,j)$ матрицы $x$.


Двумерное сечение представляет собой обычную двухиндексную


матрицу. Таким образом, у матрицы $x$ имеются двумерные сечения


ориентации $(i,j)$, $(i,t)$ и $(j,t)$. Ясно, что переставляя


местами двумерные сечения одной ориентации матрицы $x$,


представляющей собой м.\,н.\,в. многогранника $M(3,n)$, получаем новую


вершину $x'$. Эти вершины имеют одну и ту же структуру.





Число дробных компонент матрицы $x\in M(3,n)$, содержащихся в


двумерном сечении ориентации $(g,h)$ с фиксированным индексом $s$,


обозначим через $z(x^s_{gh})$. Вектор, составленный из компонент


$z(x^1_{gh}), z(x^2_{gh}), \dotsc, z(x^n_{gh})$, будем обозначать


через $z(x,(g,h))$.





Совокупность элементов матрицы $x=\|x_{ijt}\|_n$ с фиксированными


значениями двух индексов, например $i$ и $j$, будем называть


одномерным сечением ориентации $t$ матрицы $x$.





Для вершины $x=\|x_{ijt}\|_n$ многогранника $M(3,n)$ введем


множество $S(x)=\{k\in (0,1):\exists(i,j,t)\in N^3_n: 


x_{ijt}=k\}$.





Очевидно утверждение: для того чтобы пара м.\,н.\,в. $x$ и $x'$


многогранника $M(3,n)$ имела разные структуры, достаточно, чтобы


выполнялось условие $S(x)\neq S(x')$. Следующие два примера


показывают, что это условие не является необходимым.





\begin{exam}


С помощью леммы 2 и замечания 2 нетрудно убедиться, что


матрицы $x^r=\|x_{ijt}^r\|_4$, $r=0,1,2$, с ненулевыми элементами


$x_{141}^0=x_{331}^0=x_{421}^0=x_{312}^0=x_{313}^0=x_{244}^0=


x_{414}^0=x_{444}^0=\frac{1}{3}$,


$x_{122}^0=x_{233}^0=\frac{2}{3}$,


$x_{122}^1=x_{233}^1=\frac{2}{3}$,


$x_{111}^1=x_{221}^1=x_{331}^1=x_{342}^1=x_{413}^1=x_{344}^1=


x_{414}^1=x_{444}^1=\frac{1}{3}$,


$x_{122}^2=x_{233}^2=\frac{2}{3}$,


$x_{111}^2=x_{331}^2=x_{441}^2=x_{342}^2=x_{313}^2=x_{224}^2=


x_{414}^2=x_{444}^2=\frac{1}{3}$


являются м.\,н.\,в. многогранника $M(3,4)$. Вершины $x^0$, $x^1$ и


$x^2$ обладают свойствами:


\begin{itemize}


\item[1)] $S(x^0)=S(x^1)=S(x^2)=


\{\frac{1}{3},\frac{2}{3}\}$;


\item[2)] $z(x^r,(i,j))=z(x^r,(i,t))=(3,2,2,3)$, $z(x^r,(j,t))=(2,2,3,3)$,


$r=0,1,2$;


\item[3)] среди одномерных сечений матрицы $x^0$


(соответственно матрицы $x^1$ и $x^2$) имеются три (соответственно


четыре, два) сечения, каждое из которых содержит две дробных


компоненты, а остальные сечения содержат не более одной дробной


компоненты.


\end{itemize}





Из свойства 3) вытекает, что вершины $x^0$, $x^1$ и


$x^2$ имеют разные структуры.


\end{exam}





\begin{exam}


С помощью леммы 2 легко проверить, что матрицы


$x^3=\|x_{ijt}^3\|_4$ и $x^4=\|x_{ijt}^4\|_4$  с ненулевыми


элементами


$x_{111}^3=x_{221}^3=x_{242}^3=x_{414}^3=x_{444}^3=\frac{1}{4}$,


$x_{122}^3=\frac{3}{4}$,


$x_{331}^3=x_{233}^3=x_{413}^3=x_{344}^3=\frac{1}{2}$,


$x_{141}^4=x_{421}^4=x_{212}^4=x_{244}^4=x_{414}^4=\frac{1}{4}$,


$x_{122}^4=\frac{3}{4}$,


$x_{331}^4=x_{233}^4=x_{313}^4=x_{444}^4=\frac{1}{2}$ являются


м.\,н.\,в. многогранника $M(3,4)$. Вершины $x^3$ и $x^4$ обладают


свойствами:


\begin{itemize}


\item[1)] $S(x^3)=S(x^4)=\{\frac{1}{4},\frac{1}{2},\frac{3}{4}\}$;


\item[2)] $z(x^r,(i,j))=z(x^r,(i,t))=(3,2,2,3)$, $z(x^r,(j,t))=(2,3,2,3)$,


$r=3,4$;


\item[3)] среди одномерных сечений матрицы $x^3$ ($x^4$) имеются


три (два) сечения, каждое из которых содержит две дробных


компоненты, а остальные сечения содержат не более одной дробной


компоненты.


\end{itemize}





В силу свойства 3) заключаем, что вершины $x^3$ и


$x^4$ имеют разные структуры.


\end{exam}





Согласно [10] будем говорить, что м.\,н.\,в. $x$ многогранника


$M(3,n)$ имеет тип $A(p_1,q_1,p_2,q_2,p_3,q_3)$, если количество


дробных компонент, содержащихся в ее двумерных сечениях, задается


векторами


\begin{align*}


z(x,(i,j))&=(3,\dotsc,3,\underbrace{2}_{\text{место $p_1$}},


3,\dotsc,3,\underbrace{2}_{\text{место $q_1$}}


,3,\dotsc,3), \\


%


z(x,(i,t))&=(3,\dotsc,3,\underbrace{2}_{\text{место $p_2$}}


,3,\dotsc,3,\underbrace{2}_{\text{место $q_2$}},3,\dotsc,3),\\


%


z(x,(j,t))&=(3,\dotsc,3,\underbrace{2}_{\text{место $p_3$}}


,3,\dotsc,3,\underbrace{2}_{\text{место $q_3$}}


,3,\dotsc,3),


\end{align*}


где $p_l,q_l\in N_n$, $p_l\neq q_l$, $l=1,2,3$.





В [10] сформулировано предположение: значения ненулевых


элементов любой матрицы, представляющей м.\,н.\,в. многогранника


$M(3,n)$, имеющей тип $A(p_1,q_1,p_2,q_2,p_3,q_3)$, не зависят от


числа $n$ и принадлежат множеству $\{\frac{1}{3},\frac{2}{3}\}$.


Вершины $x^3$ и $x^4$ из примера 2 дают опровержение этого


предположения. Из примеров 1 и 2 вытекает также следующий


результат: тип $A(p_1,q_1,p_2,q_2,p_3,q_3)$ м.\,н.\,в. многогранника


$M(3,4)$ определяется неоднозначно, т.\,е. имеет разные структуры.





В [10] высказано также следующее предположение: для любого


натурального числа $n$ наименьшая (наибольшая) положительная


компонента среди всех м.\,н.\,в. многогранника $M(3,n)$ равна


$\frac{1}{n}$ $\left(\frac{n-1}{n}\right)$. Опровержение этого


предположения дает





\begin{thm}


Для любого числа $r\in \{0,1,2,\dotsc,n-2\}$ матрица


$x^r=\|x^r_{ijt}\|_n$ с ненулевыми


элементами


\begin{gather*}


x^r_{kk1}=\frac{1}{n+r},\ \ k=1,2,\dotsc,n-1,\\


%


x^r_{nn1}=\frac{r+1}{n+r},\ \ x^r_{k-1,k,k}=\frac{n+r-1}{n+r},


\ \ k=2,3,\dotsc,n-1,\\


%


x^r_{n-1,n,n}=\frac{n-1}{n+r},\ \ x^r_{n-1,1,k}=\frac{1}{n+r},


\ \ k=2,\dotsc,r+1\ \ (r\geq 1),\\


%


x^r_{n1k}=\frac{1}{n+r},\ \ k=r+2,\dotsc,n-1\ \ (n\geq r+3),


\ \ x^r_{n1n}=\frac{r+1}{n+r}


\end{gather*}


является м.\,н.\,в. многогранника $M(3,n)$.


\end{thm}





Эта теорема доказана в $[10]$ для случая $r=0$.





{\bf Доказательство} теоремы 1 проведем индукцией по числу $r$.


Предположим, что она верна для $r-1$.





Обозначим через $R_{ijt}$ вектор-столбец матрицы ограничений,


определяющих многогранник $M(3,n)$, т.\,е. вектор, у которого


единицы стоят в строках с номерами $i$, $(n+j)$ и $(2n+t)$, а


остальные элементы нули. В силу невырожденности вершины $x^{r-1}$


многогранника $M(3,n)$ система векторов $R_{kk1}$, $k=1,\dotsc,n$,


$R_{k-1,k,k}$, $k=2,\dotsc,n$, $R_{n-1,1,k}$, $k=2,\dotsc,r$,


$R_{n1k}$, $k=r+1,\dotsc,n$, линейно независима и вектор


$R_{n-1,1,r+1}=\sum\limits_{(i,j,t)\in K(x^r)}\beta_{ijt}R_{ijt}$, причем


$\beta_{n,1,r+1}\neq 0$. Здесь $K(x^r)=\{(i ,j,t)\in N_n^3 :


x_{ijt}^r > 0\}$. Поэтому в силу теоремы 2.1 ([17], c.\,217) система


векторов $R_{ijt}$ $\forall (i,j,t)\in K(x^r)\setminus


\{(n,1,r+1)\}$, $R_{n-1,1,r+1}$ также линейно независима. Значит,


матрица $x^r$, положительные компоненты которой отвечают линейно


независимой системе векторов, является м.н.в. многогранника


$M(3,n)$.$\quad\square$





\begin{cor}


Для любого числа $r\in \{0,1,2,\dotsc,n-2\}$ м.\,н.\,в.


$x^r$, указанная в теореме 1, обладает свойствами:


\begin{itemize}


\item[1)] $S(x^r)=\{\frac{1}{n+r},\frac{r+1}{n+r},


\frac{n-1}{n+r},\frac{n+r-1}{n+r} \}$;


\item[2)] $z(x^r,(i,j))=z(x^r,(i,t))=(n,


\underbrace{2,\dotsc,2}_{n-1})$,


$z(x^r,(j,t))=(\underbrace{2,\dotsc,2}_{n-2},r+2,n-r)$;


\item[3)] среди ее одномерных сечений имеются два сечения, одно из


которых содержит $r$, другое $n-1-r$ дробных компонент, а


остальные сечения содержат не более одной дробной компоненты.


\end{itemize}


\end{cor}





Зафиксируем число $r\in \{0,1,2,\dotsc,n-2\}$. Пусть $x^r$ --- вершина,


указанная в теореме 1. Через $\overline{x}^r$ обозначим вершину,


которая получается из вершины $x^r$ путем перестановки местами ее


двумерных сечений $n-1$ и $n$ ориентации $(j,t)$. Из следствия 2


вытекает следующий важный результат: для любого числа $r\in


\{0,1,2,\dotsc,\left[\frac{n-1}{2}\right]-1\}$ пара м.\,н.\,в.


$x^r$, $\overline{x}^{n-2-r}$ многогранника $M(3,n)$ принадлежит к


одному и тому же типу, но имеет разные структуры.





В связи с теоремой 1 возникает естественный вопрос: существуют ли


у многогранника $M(3,n)$ м.\,н.\,в., наименьшая (наибольшая)


положительная компонента которых меньше (больше) $\frac{1}{2n-2}$


$(\frac{2n-3}{2n-2})$. Ответ на этот вопрос дают теоремы 2 и 3.





\begin{thm}


У многогранника $M(3,n)$, $n\geq 6$, существует м.\,н.\,в.,


содержащая компоненту, равную $\frac{1}{n^2 -5n+6}$.


\end{thm}





\begin{pf}


Пусть $n\geq 6$, $k=[\frac{n}{2}]$. Положим


$I_1=J_1=T_1=N_k$, $I_2=J_2=T_2=N_n\setminus N_k$.





Пусть $n=2k$, $k\geq 3$. Согласно теореме 1 у многогранника


$M(I_s,J_s,T_s)$, $s=1,2$, где $|I_s|=k\geq 3$, найдутся


$(3k-2)$-нецелочисленные вершины $y$ ($y'$), содержащие по крайней


мере $2k-3$ компонент, каждая из которых равна $\frac{1}{2k-2}$


$\left(\frac{1}{2k-3}\right)$. Пусть $y^1


=\|y^1_{ijt}\|_{|I_1|\times |J_1|\times |T_1|}$ и $y^2


=\|y^2_{ijt}\|_{|I_2|\times |J_2|\times |T_2|}$ --- некоторые


$(3k-2)$-нецелочисленные вершины соответственно многогранников


$M(I_1,J_1,T_1)$ и $M(I_2,J_2,T_2)$, содержащие компоненты


$y^1_{i_1,j_1,t_1}=y^1_{i'_1,j'_1,t'_1}=\frac{1}{2k-3}$,


$y^2_{i_2,j_2,t_2}=y^2_{i'_2,j'_2,t'_2}=\frac{1}{2k-2}$, где


$(i_s,j_s,t_s),(i'_s,j'_s,t'_s)\in K(I_s,J_s,T_s,y^s)$,


$(i_s,j_s,t_s)\neq (i'_s,j'_s,t'_s)$, $s=1,2$. Тогда на основании


леммы 2 матрица $x=\|x_{ijt}\|_n$ с элементами $x_{ijt}=y^s_{ijt}$


$\forall (i,j,t)\in (I_s\times J_s\times T_s)\setminus


\{(i_s,j_s,t_s),(i'_s,j'_s,t'_s)\}$, $s=1,2$,


$x_{ijt}=\frac{1}{2k-2}$ $\forall (i,j,t)\in


\{(i_1,j_2,t_1),(i_2,j_1,t_2),(i'_2,j'_1,t'_1),(i'_1,j'_2,t'_2)\}$,


$x_{ijt}=\frac{1}{(2k-2)(2k-3)}$ $\forall (i,j,t)\in


\{(i_1,j_1,t_1),(i'_1,j'_1,t'_1)\}$, $x_{ijt}=0$ для остальных


$(i,j,t)\in N_n^3$ является м.\,н.\,в. многогранника $M(3,n)$. Так


как $n=2k$, то $\frac{1}{(2k-2)(2k-3)}=\frac{1}{n^2 -5n+6}$, что и


доказывает теорему 2 для случая $n=2k$.





Пусть теперь $n=2k+1$, $k\geq 3$. По теореме 1 у многогранника


$M(I_1,J_1,T_1)$ найдется $(3k-2)$-нецелочисленная вершина $y^1$ с


компонентами


$y^1_{i_1,j_1,t_1}=y^1_{i'_1,j'_1,t'_1}=\frac{1}{2k-2}$, а у


многогранника $M(I_2,J_2,T_2)$ (ввиду $|I_2|=|J_2|=|T_2|=k+1$) ---


$(3k+1)$-нецелочисленная вершина $y^2$ такая, что


$y^2_{i_2,j_2,t_2}=y^2_{i'_2,j'_2,t'_2}=\frac{1}{2k-1}$. Здесь


$(i_s,j_s,t_s),(i'_s,j'_s,t'_s)\in K(I_s,J_s,T_s,y^s)$,


$(i_s,j_s,t_s)\neq (i'_s,j'_s,t'_s)$, $s=1,2$. Поэтому на


основании леммы 2 у многогранника $M(3,n)$ существует м.\,н.\,в. $x$,


содержащая компоненты


$x_{i_1,j_1,t_1}=x_{i'_1,j'_1,t'_1}=\frac{1}{n^2 -5n+6}$.


\end{pf}





\begin{thm}


У многогранника $M(3,n)$, $n\geq 7$, существует м.\,н.\,в.,


содержащая компоненту, равную $\frac{n^2 -7n+10}{n^2 -7n+12}$.


\end{thm}





\begin{pf}


Пусть $n\geq 7$. Положим $I_1=J_1=T_1=N_{n-1}$,


$I_2=J_2=T_2=\{n\}$. На основании теоремы 2 у многогранника


$M(I_1,J_1,T_1)$ существует $(3|I_1|-2)$-нецелочисленная вершина


$y^1$, содержащая две компоненты, равные $\frac{1}{n^2 -7n+12}$.


Поэтому с учетом того, что многогранник $M(I_2,J_2,T_2)$


вырождается в точку $y^2_{nnn}=1$, в силу леммы 2 и замечания 2


найдется м.\,н.\,в. $x$ многогранника $M(3,n)$, содержащая компоненту,


равную $\frac{n^2 -7n+10}{n^2 -7n+12}$.


\end{pf}





С помощью теоремы 1, леммы 2 и замечания 2 аналогично доказываются


следующие три теоремы, описывающие новые типы и структуры м.\,н.\,в.


многогранника $M(3,n)$.





\begin{thm}


Для любого числа $r\in \{0,1,2,\dotsc,n-3\}$, $n\geq 4$,


матрица $x^r=\|x^r_{ijt}\|_n$ с ненулевыми


элементами


\begin{gather*}


x^r_{1n1}=x^r_{n21}=\frac{1}{n+r-1},\ \ x^r_{kk1}=\frac{1}{n+r-1},


\ \ k=3,\dotsc,n-2\ \ (n\geq 5),\\


%


x^r_{n-1,n-1,1}=\frac{r+1}{n+r-1},\ \ x^r_{k-1,k,k}=\frac{n+r-2}{n+r-1},


\ \ k=2,\dotsc,n-2,\\


%


x^r_{n-2,n-1,n-1}=\frac{n-2}{n+r-1},\ \ x^r_{n-2,1,k}=\frac{1}{n+r-1},


\ \ k=2,\dotsc,r+1\ \ (r\geq 1),\\


%


x^r_{n-1,1,k}=\frac{1}{n+r-1},\ \ k=r+2,\dotsc,n-2\ \ (n\geq r+4),\\


%


x^r_{n-1,1,n-1}=\frac{r+1}{n+r-1},\ \ x^r_{n1n}=x^r_{2nn}=\frac{1}{n+r-1},


\ \ x^r_{nnn}=\frac{n+r-3}{n+r-1}


\end{gather*}


является м.\,н.\,в. многогранника $M(3,n)$.


\end{thm}





\begin{cor}


Для любого числа $r\in \{0,1,\dotsc,n-3\}$, $n\geq 4$,


м.\,н.\,в. $x^r$, указанная в теореме 4, обладает свойствами:


\begin{itemize}


\item[1)] $S(x^r)=\{\frac{1}{n+r-1},\frac{r+1}{n+r-1},


\frac{n-2}{n+r-1},\frac{n+r-3}{n+r-1},\frac{n+r-2}{n+r-1} \}$;


\item[2)] $z(x^r,(i,j))=z(x^r,(i,t))=(n-1,\underbrace{2,2,\dotsc,2}_{n-2},3)$,


$z(x^r,(j,t))=(\underbrace{2,2,\dotsc,2}_{n-3},r+2,n-r-1,3)$;


\item[3)] среди ее одномерных сечений имеются четыре сечения, одно из


которых содержит $n-2-r$, другое $r$, третье $2$ и четвертое $2$


дробных компонент, а остальные сечения содержат не более одной


дробной компоненты.


\end{itemize}


\end{cor}





Зафиксируем $r\in \{0,1,2,\dotsc,n-3\}$. Пусть $x^r$ --- вершина,


указанная в теореме 4. Через $\overline{\overline{x}}{}^r$ обозначим


вершину, которая получается из вершины $x^r$ путем перестановки


местами ее двумерных сечений $n-2$ и $n-1$ ориентации $(j,t)$. Из


следствия 3 вытекает





\begin{cor}


Для любого числа $r\in


\{0,1,\dotsc,\left[\frac{n-2}{2}\right]-1\}$, $n\geq 4$, пара м.\,н.\,в.


$(x^r,\overline{\overline{x}}{}^{n-3-r})$ многогранника $M(3,n)$


принадлежит одному и тому же типу, но имеет разные структуры.


\end{cor}





\begin{thm}


Для любого числа $r\in \{1,2,\dotsc,n-3\}$, $n\geq 4$,


матрица $x^r=\|x^r_{ijt}\|_n$ с ненулевыми


элементами


\begin{gather*}


\displaybreak[3]


x^r_{kk1}=\frac{1}{n+r-1},\ \ k=1,\dotsc,n-2,\ \


x^r_{n-1,n-1,1}=\frac{r}{n+r-1},\\


%


x^r_{n,n-1,1}=x^r_{n-2,1,2}=\frac{1}{n+r-1},


\ \ x^r_{1n2}=\frac{n+r-2}{n+r-1},\\


%


x^r_{k-1,k,k}=\frac{n+r-2}{n+r-1},\ \ k=3,\dotsc,n-2\ \ (n\geq 5),\\


%


x^r_{n-2,n-1,n-1}=\frac{n-2}{n+r-1},


\ \ x^r_{n-2,1,k}=\frac{1}{n+r-1},\ \ k=3,\dotsc,r+1,\\


%


x^r_{n-1,1,k}=\frac{1}{n+r-1},\ \ k=r+2,\dotsc,n-2,\\


%


x^r_{n-1,1,n-1}=\frac{r+1}{n+r-1},


\ \ x^r_{n2n}=\frac{n+r-2}{n+r-1},\ \ x^r_{n-1,n,n}=\frac{1}{n+r-1}


\end{gather*}


является м.\,н.\,в. многогранника $M(3,n)$.


\end{thm}





\begin{cor}


Для любого числа $r\in \{1,2,\dotsc,n-3\}$, $n\geq 4$,


м.\,н.\,в. $x^r$, указанная в теореме 5, обладает свойствами:


\begin{itemize}


\item[1)] $S(x^r)=\{\frac{1}{n+r-1},\frac{r}{n+r-1},


\frac{r+1}{n+r-1},\frac{n-2}{n+r-1},\frac{n+r-2}{n+r-1} \}$;


\item[2)] $z(x^r,(i,j))=(n,\underbrace{2,\dotsc,2}_{n-1})$,


$z(x^r,(i,t))=(n-1,\underbrace{2,\dotsc,2}_{n-3}


,3,2)$,\\


$z(x^r,(j,t))=(\underbrace{2,\dotsc,2}_{n-3},r+2,n-r,2)$;


\item[3)] среди ее одномерных сечений имеются три сечения, одно из


которых содержит $r$, другое $n-2-r$, третье $2$ дробных


компонент, а остальные сечения содержат не более одной дробной


компоненты.


\end{itemize}


\end{cor}





Из утверждений 1) следствий 3 и 5 при $r=1$ вытекает возможность


построения различных структур м.\,н.\,в. многогранника $M(3,n)$,


$n\geq 4$, ненулевые компоненты которых определяются числами


$\frac{1}{n}$, $\frac{2}{n}$, $\frac{n-2}{n}$, $\frac{n-1}{n}$.





Зафиксируем $r\in \{1,2,\dotsc,n-3\}$, $n\geq 5$. Пусть $x^r$ ---


вершина, указанная в теореме 5. Через $\widetilde{x}^r$ обозначим


вершину, которая получается из вершины $x^r$ путем перестановки


местами ее двумерных сечений $n-2$ и $n-1$ ориентации $(j,t)$. Из


следствия 5 вытекает





\begin{cor}


Для любого числа $r\in \{1,2,\dotsc,[\frac{n-3}{2}]\}$,


$n\geq 5$, пара м.\,н.\,в. $(x^r,\widetilde{x}^{n-2-r})$ многогранника


$M(3,n)$ принадлежит к одному и тому же типу, но имеет разные


структуры.


\end{cor}





\begin{thm}


Для любого числа $r\in \{1,2,\dotsc,n-3\}$, $n\geq 4$,


матрица $x^r=\|x^r_{ijt}\|_n$ с ненулевыми


элементами


\begin{gather*}


x^r_{kk1}=\frac{1}{n+r-1},\ \ k=1,2,\dotsc,n-2,


\ \ x^r_{n-1,n,1}=\frac{r+1}{n+r-1},\\


%


x^r_{122}=\frac{r}{n+r-1},\ \ x^r_{n-2,1,2}=\frac{1}{n+r-1},


\ \ x^r_{n22}=\frac{n-2}{n+r-1},\\


%


x^r_{k-1,k,k}=\frac{n+r-2}{n+r-1},\ \ k=3,\dotsc,n-2\ \ (n\geq 5),\\


%


x^r_{n-2,n-1,n-1}=\frac{n-2}{n+r-1},


\ \ x^r_{n-2,1,k}=\frac{1}{n+r-1},\ \ k=3,\dotsc,r+1\ \ (r\geq 2),\\


%


x^r_{n-1,1,k}=\frac{1}{n+r-1},\ \ k=r+2,\dotsc,n-2,\\


%


x^r_{n-1,1,n-1}=x^r_{n,n-1,n}=\frac{r+1}{n+r-1},


\ \ x^r_{1nn}=\frac{n-2}{n+r-1}


\end{gather*}


является м.\,н.\,в. многогранника $M(3,n)$.


\end{thm}





\begin{cor}


Для любого числа $r\in \{1,2,\dotsc,n-3\}$, $n\geq 4$,


м.\,н.\,в. $x^r$, указанная в теореме 6, обладает свойствами:


\begin{itemize}


\item[1)] $S(x^r)=\{\frac{1}{n+r-1},\frac{r}{n+r-1},


\frac{r+1}{n+r-1},\frac{n-2}{n+r-1},\frac{n+r-2}{n+r-1} \}$;


\item[2)] $z(x^r,(i,j))=z(x^r,(i,t))=(n-1,3,


\underbrace{2,\dotsc,2}_{n-2})$,


$z(x^r,(j,t))=(3,\underbrace{2,\dotsc,2}_{n-4},r+2,n-r-1,2)$;


\item[3)] среди ее одномерных сечений имеются три сечения, одно из


которых содержит $n-2-r$, другое $r$, третье $2$ дробных


компонент, а остальные сечения содержат не более одной дробной


компоненты.


\end{itemize}


\end{cor}





Из следствий 3 и 7 вытекает





\begin{cor}


Для любого числа $r\in \{1,2,\dotsc,n-3\}$, $n\geq 4$,


тип м.\,н.\,в. $x^r$ многогранника $M(3,n)$, задаваемый векторами


$z(x^r,(i,j))=z(x^r,(i,t))=(n-1,\underbrace{2,\dotsc,2}_{n-2},3)$,


$z(x^r,(j,t))=(\underbrace{2,\dotsc,2}_{n-3}


,r+2,n-r-1,3)$, определяется неоднозначно.


\end{cor}





С использованием теорем 3 ([10]) и 1, а также леммы 2 и замечания


2 получены следующие две теоремы.





\begin{thm}


Матрица $x^0$ с ненулевыми элементами


\begin{gather*}


x^0_{111}=x^0_{n-3,n-3,n-3}=\frac{2}{3},


\ \ x^0_{221}=x^0_{122}=\frac{1}{12},\\


%


x^0_{2,n-2,1}=x^0_{n-1,2,2}=x^0_{n-2,2,n-2}=


x^0_{n-1,n-1,n-2}=x^0_{1,n-2,n-1}=\frac{1}{4},\\


%


x^0_{232}=x^0_{n-3,1,2}=x^0_{n-4,n-4,n-3}=\frac{1}{3},\\


%


x^0_{n,n,n-2}=x^0_{n-1,n,n}=x^0_{n,n-1,n}=\frac{1}{2},


\ \ x^0_{n-2,n-1,n-1}=\frac{3}{4},\\


%


x^0_{k-1,k-1,k}=x^0_{kkk}=x^0_{k,k+1,k}=\frac{1}{3},


\ \ k=3,\dotsc,n-4\ \ (n\geq 7),


\end{gather*}


является м.\,н.\,в. многогранника $M(3,n)$, $n\geq 6$.


\end{thm}





\begin{thm}


Матрица $x^1$ с ненулевыми элементами


\begin{gather*}


x^1_{111}=\frac{1}{6},\ x^1_{221}=x^1_{122}=


x^1_{n-3,1,2}=x^1_{n-4,n-4,n-3}=\frac{1}{3},\\


%


x^1_{nn1}=x^1_{1,1,n-2}=x^1_{n-1,n,n}=x^1_{n,n-2,n}=


\frac{1}{2},\ \ x^1_{232}=\frac{1}{12},\\


%


x^1_{2,n-1,2}=x^1_{n-2,n-2,n-2}=x^1_{n-1,n-3,n-1}=


x^1_{n-1,n-2,n-1}=\frac{1}{4},\\


%


x^1_{n-3,n-3,n-3}=\frac{2}{3},\ \ x^1_{n-2,n-1,n-1}=\frac{3}{4},\\ 


%


x^1_{k-1,k-1,k}=x^1_{kkk}=x^1_{k,k+1,k}=


\frac{1}{3},\ \ k=3,\dotsc,n-4\ \ (n\geq 7),


\end{gather*}


является м.\,н.\,в. многогранника $M(3,n)$, $n\geq 6$.


\end{thm}





Из теорем 7 и 8 вытекает





\begin{cor}


Вершины $x^0$ и $x^1$ многогранника $M(3,n)$, $n\geq


6$, обладают свойствами:


\begin{itemize}


\item[1)] $S(x^0)=\{\frac{1}{12},\frac{1}{4},\frac{1}{3},\frac{1}{2},


\frac{2}{3},\frac{3}{4} \}$, $S(x^1)=\{\frac{1}{12},\frac{1}{6},


\frac{1}{4},\frac{1}{3},\frac{1}{2},\frac{2}{3},\frac{3}{4}\}$;


\item[2)] количество дробных компонент, содержащихся в их


двумерных сечениях, задается векторами


$z(x^0,(i,j))=z(x^1,(i,j))=(3,4,2,\underbrace{3,\dotsc,3}_{n-5},2,2)$,


$z(x^0,(i,t))=(2,5,2,\underbrace{3,\dotsc,3}_


{n-5},2,2)$, $z(x^1,(i,t))=(\underbrace{3,\dotsc,3}_{n-2} ,2,2)$,


$z(x^0,(j,t))=z(x^1,(j,t))=(3,4,2,2,\underbrace{3,\dotsc,3}_{n-5} ,2)$;


\item[3)] среди их одномерных сечений имеются $2n-9$ сечений, каждое из


которых содержит две дробные компоненты, а остальные сечения


содержат не более одной дробной компоненты.


\end{itemize}


\end{cor}





\begin{thebibliography}{99}





\bibitem{1}


Balas E., Saltzman M.J. {\em Fasets of the three-index assignment


polytope} // Discrete Appl. Math. -- 1989. -- V.\,23. -- \No\,3. -- P.\,201--229.


\bibitem{2}


Pierskalla W.P. {\em The multidimensional assignment


problem} // Oper.Res. -- 1968. -- V.16. -- P.\,422--431.


\bibitem{3}


Poore A.B. {\em Multidimensional assignment formulation of data


association problems arising from multitarget and multisensor


tracking} // Comput. Optimiz. and


Applic. -- 1994. -- V.\,3. -- P.\,27--54.


\bibitem{4}


Arbib C., Pacciarelli D., Smriglio S. {\em A three-dimensional


matching model for perishable production scheduling} // Discrete


Appl. Math. -- 1999. -- \No\,92. -- P.\,1--15.


\bibitem{5}


Euler R. {\em Odd cycles and a class of facets of the axial


$3$-index assignment polytope} // Zastosowania


Matematyki. -- 1987. -- V.19. -- \No\,3--4. -- P.\,375--386.


\bibitem{6}


Gwan G., Qi L. {\em On facets of the three-index assignment


polytope} // Australasian J. Combinatorics. -- 1992. -- V.\,6. -- P.\,67--87.


\bibitem{7}


Qi L., Balas E., Gwan G. {\em A new facet class and a polyhedral


method for the three-index assignment problem} // In: D.\,Du and


J.\,Sun (Eds) Advances in Optimization and Approximation. -- Kluwer


Academic. -- 1994. -- P.\,256--274.


\bibitem{8}


Кравцов М.К., Лукшин Е.В. {\em О нецелочисленных вершинах


многогранника многоиндексной аксиальной задачи выбора} // Изв. вузов.


Математика. -- 1999. -- \No\,12. -- C.\,65--70.


\bibitem{9}


Кравцов М.К., Кравцов В.М., Лукшин Е.В. {\em О числе


$r$-нецелочисленных вершин многогранника трехиндексной аксиальной


задачи выбора} // Изв. вузов. Mатематика. -- 2000. -- \No\,12. -- C.\,89--92.


\bibitem{10}


Кравцов М.К., Кравцов В.М., Лукшин Е.В. {\em О типах


$(3n-2)$-нецелочисленных вершин многогранника трехиндексной


аксиальной задачи выбора} // Изв. вузов.


Mатематика. -- 2002. -- \No\,12. -- С.\,84--90.


\bibitem{11}


Кравцов М.К., Кравцов В.М., Лукшин Е.В. {\em О нецелочисленных


вершинах многогранника трехиндексной аксиальной задачи о


назначениях} // Дискретн. матем. -- 2001. -- Т.\,13. -- Вып.\,2. -- С.\,120--143.


\bibitem{12}


Кравцов М.К., Кравцов В.М., Лукшин Е.В. {\em О числе


нецелочисленных вершин многогранника трехиндексной аксиальной


задачи о назначениях} // Весцi НАН Беларусi. Сер. фiз.-мат.


навук. -- 2000. -- \No\,4. -- С.\,59--65.


\bibitem{13}


Кравцов В.М. Оценки снизу числа нецелочисленных вершин


многогранника трехиндексной аксиальной задачи о


назначениях//Вестник БГУ. Сер.1.--2002.--N3.--С.87--90.


\bibitem{14}


Кравцов В.М. {\em О новых типах максимально нецелочисленных вершин


многогранника трехиндексной аксиальной задачи о


назначениях} // Вестник БГУ. Сер.1. -- 2003. -- \No\,3. -- С.\,80--85.


\bibitem{15}


Емеличев В.А., Ковалев М.М., Кравцов М.К. {\em Многогранники,


графы, оптимизация}. -- М.: Наука, 1981. -- 342~с.


\bibitem{16}


Емеличев В.А., Кравцов М.К. {\em Полиэдральные аспекты


многоиндексных аксиальных транспортных задач} // Дискретн.


матем. -- 1991. -- Т.\,3. -- Вып.\,2. -- С.\,3--24.


\bibitem{17}


Юдин Д.Б., Гольштейн Е.Г. {\em Линейное программирование. Теория


и конечныe методы}. -- М.: Физико-математическая литература, 1963. -- 776~с.





\end{thebibliography}





\vskip 0.5cm





\post{\it Белорусский государственный}{\it Поступила}


\post{\it университет}{05.05.2004}





\label{end}


\end{document}


